\section{图表与数字}
\label{sec:figure}

表格一律用 \texttt{booktabs} 的三线表：顶线、中线、底线，不加竖线。表题在表
上方，图题在图下方，两者都用竖线把标签和说明分开（\texttt{表 1 | ...}）。
最优值加粗（\verb|\best{}|），次优下划线（\verb|\second{}|），并在题注里
写明含义，读者不必回头翻正文。

每张图表都要在正文里被引用，否则删掉它。中文里章节用 \verb|\secref{}|，
写成「第 3 节」；图表用 \verb|\cref{}|，写成「图 1」「表 2」。这样改章节顺序
时不会留下错误的「第 3 节」。

图~\ref{fig:arch} 是仓库结构与构建链路，用 TikZ 直接画在文档里，不依赖外部
图片，改结构时改代码即可。

%% 仓库结构与构建链路示意图（TikZ，无外部图片依赖）
\begin{figure}[htbp]
  \centering
  \begin{tikzpicture}[
      font=\small,
      box/.style={draw=bgrule, rounded corners=2pt, fill=bgbrand!4,
                  align=center, inner sep=4pt, minimum height=8mm,
                  minimum width=40mm},
      core/.style={draw=bgbranddeep, rounded corners=3pt, fill=bgbrand!12,
                   align=center, inner sep=6pt, minimum height=26mm,
                   minimum width=38mm},
      outbox/.style={draw=bgrule, rounded corners=2pt, fill=black!4,
                  align=center, inner sep=4pt, minimum height=9mm,
                  minimum width=34mm},
      arr/.style={-{Latex[length=2mm]}, bgrule!90, thick},
      note/.style={font=\scriptsize, text=bgink, align=center},
      node distance=4mm and 16mm,
    ]
    \node[box] (mains) {main-zh.tex / main-en.tex};
    \node[box, below=4mm of mains] (secs) {sections/};
    \node[box, below=4mm of secs] (figs) {figures/};
    \node[box, below=4mm of figs] (bib) {refs.bib};

    \node[core, right=18mm of secs, yshift=-4mm] (cls)
      {\textbf{boogonreport.cls}\\[2pt]{\scriptsize 文档类}};

    \node[outbox, right=16mm of cls, yshift=12mm] (pdf) {main-zh.pdf / main-en.pdf};
    \node[outbox, right=16mm of cls, yshift=-12mm] (ci) {GitHub Actions};

    \draw[arr] (mains.east) -- (cls.west |- mains.east);
    \draw[arr] (bib.east) -- (cls.west |- bib.east);
    \draw[arr] (cls.east) -- ++(6mm,0) |- (pdf.west);
    \draw[arr] (cls.east) -- ++(6mm,0) |- (ci.west);

    \node[note, below=6mm of cls] {构建：build.ps1 或 latexmk\\英文 pdflatex，中文 xelatex};
  \end{tikzpicture}
  \caption{仓库结构与构建链路。内容与类文件分离：换项目时只复制
  \texttt{boogonreport.cls} 与自己的 \texttt{sections/}、\texttt{figures/}，
  构建脚本与 CI 保持不变。}
  \label{fig:arch}
\end{figure}


数字的写法比字体更容易出问题，因为错误不会报错，只会让读者困惑。表~\ref{tab:numbers}
是我们在项目里实际踩过的几处。

\begin{table}[htbp]
  \centering
  \small
  \caption{数字写法约定。中间是排版后的样子，右边是应当写进源码的形式。}
  \label{tab:numbers}
  \begin{tabular}{@{}lll@{}}
    \toprule
    场景 & 排版结果 & 源码写法 \\
    \midrule
    千分位        & 32,768      & \texttt{32\{,\}768} \\
    乘号          & 8$\times$H100 & \texttt{8\$\textbackslash times\$H100} \\
    数值与单位    & 38\,ms      & \texttt{38\textbackslash,ms} \\
    百分号        & 62.4\%      & \texttt{62.4\textbackslash\%} \\
    精度统一      & 62.4 与 61.0 & 同一张表不要混用 62.40 \\
    提升幅度      & +2.1 points over the 7B baseline & 写清基线，不要只写 +2.1 \\
    \bottomrule
  \end{tabular}
\end{table}

第三行常被忽略：\texttt{38ms} 在断行时可能被拆成两行，读起来像 38 和 ms 两个
东西；\verb|38\,ms| 用一个细空格把数值和单位绑在一起，断行时不会被拆开。
